\documentclass{beamer}

\batchmode

\usepackage{amsmath,amssymb,enumerate,epsfig,bbm,calc,color,ifthen,capt-of}

\usetheme{Berlin}
\usecolortheme{mit}

%% <language_settings> %%
\usepackage[utf8x]{inputenc}
\usepackage[czech]{babel}
%% </language_settings> %%

%% Source code insertions
\usepackage{verbatim}

%% <change me> %%
\title{Gentoo KDE: stable, fresh, and bleeding edge!}
\author[Tomáš Chvátal]{Tomáš Chvátal $<$scarabeus@gentoo.org$>$}
\date{2012/10/20}
%% </change me> %%

%% <logo> %%
\usepackage[absolute,overlay]{textpos}
\setlength{\TPHorizModule}{1mm}
\setlength{\TPVertModule}{1mm}
\newcommand{\MyLogo}{
	\begin{textblock}{14}(110,7)
		\includegraphics[width=14mm,height=14mm]{gentoo-logo.png}
	\end{textblock}
}

\newcommand*\oldmacro{}
\let\oldmacro\insertshorttitle
\renewcommand*\insertshorttitle{
	\MyLogo
	\oldmacro\hfill
}
%% </logo> %%

%% <bgimage> %%
\usebackgroundtemplate{
	\includegraphics[width=\paperwidth,height=\paperheight]{gentoo-background-1.png}
}
%% </bgimage> %%

%% <ToC Before each section> %%
%\AtBeginSection[]{
%	\begin{frame}<beamer>
%		\frametitle{Přehled}
%		\tableofcontents[currentsection]
%	\end{frame}
%}
%% </ToC Before each section> %%

%% <each bullet on one slide> %%
%% better to use \pause on points where we really want to stop :)
%\beamerdefaultoverlayspecification{<+->}
%% </each bullet on one slide> %%
% -----------------------------------------------------------------------------
\begin{document}
% -----------------------------------------------------------------------------
\frame{\titlepage}
%\section[Přehled]{}
%\begin{frame}{Přehled}
%	\tableofcontents
%\end{frame}
% -----------------------------------------------------------------------------
\section{Introduction}

\subsection{Tomáš}

\begin{frame}{Who the hell is Tomáš Chvátal}
	\begin{itemize}
		\item Council member (currently second term)
		\item KDE team member (former lead)
		\item Libreoffice maintainer
		\item Formely also working on X11, Overlays, Clustering, QA, ...
	\end{itemize}
\end{frame}

% -----------------------------------------------------------------------------
\section{What is KDE}

\subsection{Overview}

\begin{frame}{What is KDE}
	\begin{itemize}
		\item Fancy Qt based desktop environment with many features.
		\item Visually appealing and configurable.
		\item Bloated thing wrt database usage.
	\end{itemize}
\end{frame}

\begin{frame}{How to make not pain and really usable}
	\begin{itemize}
		\item Make the nepomuk (aka semantic-desktop) optional for those who prefer to use their system rather than index its content all the time.
		\item Make akonadi also optional and allow installation of pre-akonadi versions to have something really usable :-)
		\item Make all this fun run on hardened and so make it bit safer.
	\end{itemize}
\end{frame}

\begin{frame}{Running KDE on gentoo}
	\begin{itemize}
		\item Use desktop or desktop/kde profile
		\item Enable devicekit and consolekit useflag
		\item Select the packages that are actually needed
	\end{itemize}
\end{frame}

\begin{frame}{What packages would I need}
	\begin{itemize}
		\item kde-base/kde-meta: Not really, this is metapackage that includes EVERYTHING.
		\item kde-base/kdebase-meta and kde-base/kdebase-runtime-meta: Basic subset on which one can install additional packages.
		\item kde-base/*-meta: These are joint packages per various category.
	\end{itemize}
\end{frame}

\begin{frame}{Interesting useflags}
	\begin{itemize}
		\item semantic-desktop: useflag removing all akonadi, virtuoso and nepomuk dependencies.
		\item udisks/upower: enable or disable dependency over these tools, it is runtime detected tho.
	\end{itemize}
\end{frame}

\subsection{Bleedingness}
\begin{frame}{Stable}
	\begin{itemize}
		\item Bit older than latest released stable version from upstream, but tested a lot to ensure its stability.
		\item Usually new stable version is done every other month.
	\end{itemize}
\end{frame}

\begin{frame}{Testing}
	\begin{itemize}
		\item Always 0 day version bump with latest upstream released stuff.
		\item Usually contains a lots of issues that upstream didn't care enough to fix.
		\item This stuff is taken from overlay by fancy update script in bash I wrote 4 years ago.
	\end{itemize}
\end{frame}

\begin{frame}{Overlay - aka snapshots / live stuff ; let me break your boxen TM}
	\begin{itemize}
		\item Usually live versions of master branch and currently produced branch (ie. 4.8 when 4.9 is in development).
		\item The overlay is the place where all the development is done.
		\item The branch ebuilds are usually just the bug fixes to next release.
		\item Live ebuilds is where magic happens, and you will usually cry for help when stuff breaks.
	\end{itemize}
\end{frame}

\section{What can I do for KDE in Gentoo}

\subsection{Possible contributions}

\begin{frame}{Contributing}
	\begin{itemize}
		\item Join \#gentoo-kde on freenode.net and chat with us about your secret plans.
		\item Clone up the git repository and help others by using live ebuilds and fix the issues.
		\item Fire up bugzilla and help us closing issues (we like to have just few bugs, something around 100 is perfect).
	\end{itemize}
\end{frame}

\begin{frame}{KDE stable}
	\begin{itemize}
		\item New KDE subproject with focus on stable users.
		\item Team simply works on testing various packages from the testing tree and moves them for stabilisation if they meet the requirements.
		\item Actually anyone who is using stable system (nothing much keyworded) can join this team and help up.
		\item Planned results for this team are to provide even more polished KDE experience on stable systems.
	\end{itemize}
\end{frame}

\section{KDE SC 5}
\begin{frame}{Few ideas what to do for KDE 5 release}
	\begin{center}Bit of discussion about KDE future.\end{center}
\end{frame}

% -----------------------------------------------------------------------------
\section{Wrap-up}

\subsection{Wrap-up}

\begin{frame}{Links to learn from and links to stuff I mentioned}
	\begin{itemize}
		\item http://wiki.gentoo.org/wiki/KDE
		\item http://www.gentoo.org/proj/en/desktop/kde/
		\item http://git.overlays.gentoo.org/gitweb/?p=proj/kde.git
	\end{itemize}
\end{frame}

\begin{frame}{Questions}
	\begin{center}Nobody wants to ask anything. Right?\end{center}
\end{frame}

\begin{frame}{Thank you}
	\begin{center}Thank you for your attention!\end{center}
\end{frame}
% -----------------------------------------------------------------------------
\end{document}
